\chapter{Smooth Maps and Diffeomorphisms}
\label{ch:vii03-smooth-maps-and-diffeomorphisms}

With smooth structures fixed, a map is smooth when its Euclidean coordinate representatives are smooth:
\[
\boxed{F\text{ smooth}\iff \psi\circ F\circ\varphi^{-1}\text{ smooth locally}.}
\]

\section{Smooth maps}
\begin{definition}[Smooth map at a point]\label{def:vii03-smooth-map}
\(F:M\to N\) is smooth at \(p\) if some compatible source and target charts make \(\psi\circ F\circ\varphi^{-1}\) smooth near \(\varphi(p)\).
\end{definition}
\begin{theorem}[Chart independence]\label{thm:vii03-chart-independence}
Smoothness in one compatible chart pair implies smoothness in every compatible chart pair.
\end{theorem}
\begin{proof}
New coordinate representatives differ by pre- and post-composition with smooth transition maps.
\end{proof}
\begin{definition}[Smooth map]\label{def:vii03-global-smooth}
A map is smooth if it is smooth at every point.
\end{definition}

\section{Local criteria}
\begin{theorem}[Locality]\label{thm:vii03-locality}
Smoothness can be checked on an arbitrary open cover of the source.
\end{theorem}
\begin{proposition}[Componentwise criterion]\label{prop:vii03-componentwise}
A map \(M\to\mathbb R^m\) is smooth iff all its components are smooth real-valued functions.
\end{proposition}

\section{Identity and composition}
\begin{proposition}[Identity]\label{prop:vii03-identity}
Identity maps are smooth.
\end{proposition}
\begin{theorem}[Composition]\label{thm:vii03-composition}
Compositions of smooth maps are smooth.
\end{theorem}

\section{Diffeomorphisms}
\begin{definition}[Diffeomorphism]\label{def:vii03-diffeomorphism}
A diffeomorphism is a bijective smooth map with smooth inverse.
\end{definition}
\begin{warning}[Smooth bijection is insufficient]\label{warning:vii03-bijection}
\(x\mapsto x^3\) is smooth and bijective on \(\mathbb R\), but its inverse is not smooth at zero.
\end{warning}

\section{An explicit model diffeomorphism}
\begin{theorem}[Ball to Euclidean space]\label{thm:vii03-ball}
For \(B_r\subset\mathbb R^k\),
\[
F(x)=\frac{rx}{\sqrt{r^2-\|x\|^2}}
\]
is a diffeomorphism onto \(\mathbb R^k\), with inverse
\[
G(y)=\frac{ry}{\sqrt{r^2+\|y\|^2}}.
\]
\end{theorem}
\begin{corollary}[Full-space charts]\label{cor:vii03-full-space}
Every point of a smooth \(k\)-manifold has a chart onto all of \(\mathbb R^k\).
\end{corollary}

\section{Dimension invariance}
\begin{theorem}[Euclidean dimension obstruction]\label{thm:vii03-dimension}
Nonempty Euclidean opens can be diffeomorphic only when their dimensions agree.
\end{theorem}
\begin{corollary}[Manifold dimension]\label{cor:vii03-manifold-dimension}
Diffeomorphic smooth manifolds have the same local dimension.
\end{corollary}

\section{Category viewpoint}
\begin{remark}[Smooth category]\label{rem:vii03-category}
Smooth manifolds and smooth maps form a category; its isomorphisms are diffeomorphisms.
\end{remark}
\section{Exercises}
\begin{exercise}\label{exr:vii03-01}
State smoothness of \(F:M\to N\) at \(p\) in coordinates.
\end{exercise}
\begin{hint}
% PEDAGOGY-ENRICHED-VII
Choose source and target charts. Make the controlling geometric criterion explicit before the calculation. Move the map into local coordinates, differentiate there, and then check that the conclusion is independent of the chosen charts.
\end{hint}
\begin{solution}
Require \(\psi\circ F\circ\varphi^{-1}\) to be smooth near \(\varphi(p)\).
\end{solution}

\begin{exercise}\label{exr:vii03-02}
Prove chart-independence of smoothness.
\end{exercise}
\begin{hint}
% PEDAGOGY-ENRICHED-VII
Factor through transition maps. Work in a convenient local model first, then return to the invariant statement. Move the map into local coordinates, differentiate there, and then check that the conclusion is independent of the chosen charts.
\end{hint}
\begin{solution}
Any new representative equals a smooth target transition composed with the old representative and a smooth source transition.
\end{solution}

\begin{exercise}\label{exr:vii03-03}
Show the identity map is smooth.
\end{exercise}
\begin{hint}
% PEDAGOGY-ENRICHED-VII
It is a transition map in coordinates. After the main computation, check the hypotheses and geometric interpretation. Move the map into local coordinates, differentiate there, and then check that the conclusion is independent of the chosen charts.
\end{hint}
\begin{solution}
Its coordinate representation is \(\psi\circ\varphi^{-1}\), smooth by compatibility.
\end{solution}

\begin{exercise}\label{exr:vii03-04}
Show constant maps are smooth.
\end{exercise}
\begin{hint}
% PEDAGOGY-ENRICHED-VII
Coordinate representatives are constant. Make the controlling geometric criterion explicit before the calculation. Move the map into local coordinates, differentiate there, and then check that the conclusion is independent of the chosen charts.
\end{hint}
\begin{solution}
Constant Euclidean maps are smooth.
\end{solution}

\begin{exercise}\label{exr:vii03-05}
Prove compositions of smooth maps are smooth.
\end{exercise}
\begin{hint}
% PEDAGOGY-ENRICHED-VII
Choose an intermediate chart. Work in a convenient local model first, then return to the invariant statement. Move the map into local coordinates, differentiate there, and then check that the conclusion is independent of the chosen charts.
\end{hint}
\begin{solution}
The coordinate representative of a composite is the composite of Euclidean representatives.
\end{solution}

\begin{exercise}\label{exr:vii03-06}
Prove smoothness is local on an open cover.
\end{exercise}
\begin{hint}
% PEDAGOGY-ENRICHED-VII
Check each point in one cover member. After the main computation, check the hypotheses and geometric interpretation. Move the map into local coordinates, differentiate there, and then check that the conclusion is independent of the chosen charts.
\end{hint}
\begin{solution}
Every point has a neighborhood on which the restriction is smooth, hence the map is smooth at every point.
\end{solution}

\begin{exercise}\label{exr:vii03-07}
Prove \(F:M\to\mathbb R^k\) is smooth iff each component is smooth.
\end{exercise}
\begin{hint}
% PEDAGOGY-ENRICHED-VII
Use the Euclidean componentwise criterion. Make the controlling geometric criterion explicit before the calculation. Move the map into local coordinates, differentiate there, and then check that the conclusion is independent of the chosen charts.
\end{hint}
\begin{solution}
In source coordinates the representative is vector-valued and is smooth exactly when all components are smooth.
\end{solution}

\begin{exercise}\label{exr:vii03-08}
Relate VII/02 smooth functions to smooth maps \(M\to\mathbb R\).
\end{exercise}
\begin{hint}
% PEDAGOGY-ENRICHED-VII
Use the identity target chart. Work in a convenient local model first, then return to the invariant statement. Move the map into local coordinates, differentiate there, and then check that the conclusion is independent of the chosen charts.
\end{hint}
\begin{solution}
The two coordinate definitions are identical.
\end{solution}

\begin{exercise}\label{exr:vii03-09}
Define diffeomorphism and prove it is an equivalence relation.
\end{exercise}
\begin{hint}
% PEDAGOGY-ENRICHED-VII
Use identity, inverse, composition. After the main computation, check the hypotheses and geometric interpretation. Move the map into local coordinates, differentiate there, and then check that the conclusion is independent of the chosen charts.
\end{hint}
\begin{solution}
All three properties follow from smooth identity maps and closure of smooth maps under composition.
\end{solution}

\begin{exercise}\label{exr:vii03-10}
Give a smooth bijection that is not a diffeomorphism.
\end{exercise}
\begin{hint}
% PEDAGOGY-ENRICHED-VII
Use \(x^3\). Make the controlling geometric criterion explicit before the calculation. Move the map into local coordinates, differentiate there, and then check that the conclusion is independent of the chosen charts.
\end{hint}
\begin{solution}
\(x\mapsto x^3\) is smooth and bijective, but \(x^{1/3}\) is not smooth at zero.
\end{solution}

\begin{exercise}\label{exr:vii03-11}
Show affine isomorphisms of \(\mathbb R^n\) are diffeomorphisms.
\end{exercise}
\begin{hint}
% PEDAGOGY-ENRICHED-VII
Write the affine inverse. Work in a convenient local model first, then return to the invariant statement. Move the map into local coordinates, differentiate there, and then check that the conclusion is independent of the chosen charts.
\end{hint}
\begin{solution}
\(x\mapsto Ax+b\) has smooth inverse \(y\mapsto A^{-1}(y-b)\).
\end{solution}

\begin{exercise}\label{exr:vii03-12}
Show restrictions of smooth maps to open subsets are smooth.
\end{exercise}
\begin{hint}
% PEDAGOGY-ENRICHED-VII
Restrict coordinate representatives. After the main computation, check the hypotheses and geometric interpretation. Move the map into local coordinates, differentiate there, and then check that the conclusion is independent of the chosen charts.
\end{hint}
\begin{solution}
Restrictions of smooth Euclidean maps are smooth.
\end{solution}

\begin{exercise}\label{exr:vii03-13}
Show corestriction to an open submanifold is smooth.
\end{exercise}
\begin{hint}
% PEDAGOGY-ENRICHED-VII
Use induced target charts. Make the controlling geometric criterion explicit before the calculation. Move the map into local coordinates, differentiate there, and then check that the conclusion is independent of the chosen charts.
\end{hint}
\begin{solution}
Target charts are restrictions of ambient charts, giving the same coordinate expressions.
\end{solution}

\begin{exercise}\label{exr:vii03-14}
Verify \(F(x)=rx/\sqrt{r^2-\|x\|^2}\) is smooth on \(B_r\).
\end{exercise}
\begin{hint}
% PEDAGOGY-ENRICHED-VII
The denominator is positive. Work in a convenient local model first, then return to the invariant statement. Move the map into local coordinates, differentiate there, and then check that the conclusion is independent of the chosen charts.
\end{hint}
\begin{solution}
It is built from polynomial, square-root and reciprocal operations on a positive function.
\end{solution}

\begin{exercise}\label{exr:vii03-15}
Find the inverse of the ball map.
\end{exercise}
\begin{hint}
% PEDAGOGY-ENRICHED-VII
Solve radially. After the main computation, check the hypotheses and geometric interpretation. Move the map into local coordinates, differentiate there, and then check that the conclusion is independent of the chosen charts.
\end{hint}
\begin{solution}
\(G(y)=ry/\sqrt{r^2+\|y\|^2}\).
\end{solution}

\begin{exercise}\label{exr:vii03-16}
Prove that inverse is smooth everywhere.
\end{exercise}
\begin{hint}
% PEDAGOGY-ENRICHED-VII
The denominator never vanishes. Make the controlling geometric criterion explicit before the calculation. Move the map into local coordinates, differentiate there, and then check that the conclusion is independent of the chosen charts.
\end{hint}
\begin{solution}
\(r^2+\|y\|^2>0\) for all \(y\).
\end{solution}

\begin{exercise}\label{exr:vii03-17}
Use the ball diffeomorphism to get a chart onto all \(\mathbb R^n\).
\end{exercise}
\begin{hint}
% PEDAGOGY-ENRICHED-VII
Shrink to a coordinate ball. Work in a convenient local model first, then return to the invariant statement. Move the map into local coordinates, differentiate there, and then check that the conclusion is independent of the chosen charts.
\end{hint}
\begin{solution}
Restrict a chart to the preimage of a ball and compose with the ball-to-space diffeomorphism.
\end{solution}

\begin{exercise}\label{exr:vii03-18}
Prove diffeomorphic Euclidean opens have equal dimension.
\end{exercise}
\begin{hint}
% PEDAGOGY-ENRICHED-VII
Differentiate inverse identities. After the main computation, check the hypotheses and geometric interpretation. Move the map into local coordinates, differentiate there, and then check that the conclusion is independent of the chosen charts.
\end{hint}
\begin{solution}
The derivative of a diffeomorphism has a two-sided linear inverse, forcing equal dimensions.
\end{solution}

\begin{exercise}\label{exr:vii03-19}
Show every diffeomorphism is a homeomorphism.
\end{exercise}
\begin{hint}
% PEDAGOGY-ENRICHED-VII
Smooth maps are continuous. Make the controlling geometric criterion explicit before the calculation. Move the map into local coordinates, differentiate there, and then check that the conclusion is independent of the chosen charts.
\end{hint}
\begin{solution}
A diffeomorphism and its smooth inverse are continuous inverse maps.
\end{solution}

\begin{exercise}\label{exr:vii03-20}
Show changing the target chart preserves smoothness.
\end{exercise}
\begin{hint}
% PEDAGOGY-ENRICHED-VII
Post-compose by transition. Work in a convenient local model first, then return to the invariant statement. Move the map into local coordinates, differentiate there, and then check that the conclusion is independent of the chosen charts.
\end{hint}
\begin{solution}
The new representative is a smooth transition composed with the old representative.
\end{solution}

\begin{exercise}\label{exr:vii03-21}
Show changing the source chart preserves smoothness.
\end{exercise}
\begin{hint}
% PEDAGOGY-ENRICHED-VII
Pre-compose by transition. After the main computation, check the hypotheses and geometric interpretation. Move the map into local coordinates, differentiate there, and then check that the conclusion is independent of the chosen charts.
\end{hint}
\begin{solution}
The new representative is the old one composed with a smooth source transition.
\end{solution}

\begin{exercise}\label{exr:vii03-22}
Explain why a generating atlas suffices to check smoothness.
\end{exercise}
\begin{hint}
% PEDAGOGY-ENRICHED-VII
All maximal charts are compatible with it. Make the controlling geometric criterion explicit before the calculation. Move the map into local coordinates, differentiate there, and then check that the conclusion is independent of the chosen charts.
\end{hint}
\begin{solution}
Chart-independence propagates smoothness from generating charts to every chart in the maximal atlas.
\end{solution}

\begin{exercise}\label{exr:vii03-23}
Explain why a homeomorphism need not be a diffeomorphism.
\end{exercise}
\begin{hint}
% PEDAGOGY-ENRICHED-VII
Continuity is weaker than differentiability. Work in a convenient local model first, then return to the invariant statement. Move the map into local coordinates, differentiate there, and then check that the conclusion is independent of the chosen charts.
\end{hint}
\begin{solution}
The \(x^3\) example is a homeomorphism whose inverse is not smooth.
\end{solution}

\begin{exercise}\label{exr:vii03-24}
State the bridge to tangent spaces.
\end{exercise}
\begin{hint}
% PEDAGOGY-ENRICHED-VII
Smooth maps have first-order linearizations. After the main computation, check the hypotheses and geometric interpretation. Move the map into local coordinates, differentiate there, and then check that the conclusion is independent of the chosen charts.
\end{hint}
\begin{solution}
VII/04 defines \(dF_p:T_pM\to T_{F(p)}N\), the coordinate-free derivative.
\end{solution}

\section{Solved problem dossiers}
\begin{problem}[Legacy Problem 3: local chart criterion]\label{prob:vii03-01}
Prove it suffices to cover \(M\) by open sets on which one chosen coordinate representative of \(F\) is smooth.
\end{problem}
\begin{solution}
At each point use the representative supplied by its cover member; chart-independence upgrades it to all compatible coordinates.
\end{solution}

\begin{problem}[Legacy diffeomorphism problem: ball to space]\label{prob:vii03-02}
Prove \(F(x)=rx/\sqrt{r^2-\|x\|^2}\) is a diffeomorphism \(B_r\to\mathbb R^k\).
\end{problem}
\begin{solution}
The inverse is \(G(y)=ry/\sqrt{r^2+\|y\|^2}\); both are smooth and direct substitution gives both identities.
\end{solution}

\begin{problem}[Legacy consequence: full-space charts]\label{prob:vii03-03}
Show every point of a smooth \(k\)-manifold admits a chart onto all of \(\mathbb R^k\).
\end{problem}
\begin{solution}
Shrink any chart to a coordinate ball and compose with the legacy ball diffeomorphism.
\end{solution}

\begin{problem}[Legacy dimension obstruction]\label{prob:vii03-04}
Prove nonempty opens in \(\mathbb R^n\) and \(\mathbb R^m\) cannot be diffeomorphic when \(n\neq m\).
\end{problem}
\begin{solution}
Differentiate the inverse identities; the derivative must be a linear isomorphism.
\end{solution}

\begin{problem}[Chart-independence theorem]\label{prob:vii03-05}
Prove the definition of smooth map is independent of both source and target charts.
\end{problem}
\begin{solution}
Factor a new representative through the old representative and two smooth transition maps.
\end{solution}

\begin{problem}[Composition theorem]\label{prob:vii03-06}
Prove composition of smooth manifold maps is smooth.
\end{problem}
\begin{solution}
Choose compatible source, middle and target charts so the coordinate representative is a Euclidean composition.
\end{solution}

\begin{problem}[Smooth bijection warning]\label{prob:vii03-07}
Give a smooth bijection with nonsmooth inverse.
\end{problem}
\begin{solution}
Use \(x\mapsto x^3\), whose inverse \(x^{1/3}\) is not differentiable at zero.
\end{solution}

\begin{problem}[Identity between two smooth structures]\label{prob:vii03-08}
Characterize when the identity homeomorphism between two smooth structures on one topological manifold is a diffeomorphism.
\end{problem}
\begin{solution}
Exactly when the structures are equivalent, because its coordinate representatives are the cross-transition maps.
\end{solution}

\begin{problem}[Componentwise criterion]\label{prob:vii03-09}
Prove \(F:M\to\mathbb R^m\) is smooth iff all coordinate components are smooth functions.
\end{problem}
\begin{solution}
Apply the Euclidean componentwise smoothness theorem in any source chart.
\end{solution}

\begin{problem}[Open-cover locality]\label{prob:vii03-10}
If \(M=\bigcup U_\alpha\) and every \(F|_{U_\alpha}\) is smooth, prove \(F\) smooth.
\end{problem}
\begin{solution}
Smoothness at a point is witnessed in any cover member containing that point.
\end{solution}

\begin{problem}[Dimension invariance]\label{prob:vii03-11}
Show diffeomorphic connected smooth manifolds have equal dimension.
\end{problem}
\begin{solution}
Choose charts at corresponding points; the induced Euclidean diffeomorphism forces equal chart dimensions.
\end{solution}

\begin{problem}[Category viewpoint]\label{prob:vii03-12}
Show smooth manifolds and smooth maps form a category and identify its isomorphisms.
\end{problem}
\begin{solution}
Identity and composition are smooth; categorical isomorphisms are exactly diffeomorphisms.
\end{solution}

\section{Challenges}
\begin{problem}[Challenge]\label{chal:vii03-01}
Give a smooth homeomorphism that is not a diffeomorphism and diagnose failure of the inverse.
\end{problem}
\begin{solution}
\(x\mapsto x^3\); the inverse cube root is not differentiable at zero.
\end{solution}

\begin{problem}[Challenge]\label{chal:vii03-02}
Prove \(F:M\to N\) is smooth iff every local smooth real-valued \(g\) on \(N\) pulls back to a smooth \(g\circ F\).
\end{problem}
\begin{solution}
The forward direction is composition. For the converse choose target coordinate functions; their pullbacks are the components of the coordinate representative.
\end{solution}

\begin{problem}[Challenge]\label{chal:vii03-03}
Analyze the radial behavior of the ball diffeomorphism near \(\partial B_r\).
\end{problem}
\begin{solution}
For \(x=\rho u\), the radial factor \(r\rho/\sqrt{r^2-\rho^2}\) tends to \(+\infty\) as \(\rho\to r^{-}\).
\end{solution}

\begin{problem}[Challenge]\label{chal:vii03-04}
Prove a gluing principle for smooth maps on an open cover.
\end{problem}
\begin{solution}
Topological gluing gives the map; smoothness follows locally because the glued map agrees with a smooth piece near every point.
\end{solution}

\begin{problem}[Challenge]\label{chal:vii03-05}
Design a practical diffeomorphism checklist.
\end{problem}
\begin{solution}
Prove bijectivity, smoothness locally in convenient charts, and smoothness of the inverse; do not infer inverse smoothness from bijectivity.
\end{solution}

\section*{Bridge to VII/04}
VII/04 constructs tangent spaces and the differential \(dF_p\), turning the first derivative of a coordinate representative into an intrinsic linear map.
